%!TEX program = xelatex

\documentclass[11pt,a4paper]{extarticle}

% ---------------------------------------------------------------------
%  LIBRERÍA Y FUENTES DESDE GITHUB (siempre la última versión)
%  https://github.com/nataliagamezbarea/apuntes-latex
%
%  Al compilar: descarga estilo.tex + fonts/ de GitHub a una carpeta
%  temporal, los usa, y al terminar la BORRA (no queda nada, ni en el
%  proyecto ni en tu equipo). Si algo cambia en GitHub, la próxima
%  compilación lo coge.
%  Necesita -shell-escape (ya está activado en la config de VS Code).
% ---------------------------------------------------------------------
\immediate\write18{rm -rf "/tmp/apuntes-latex-lib-\jobname" && mkdir -p "/tmp/apuntes-latex-lib-\jobname" && curl -fsSL https://codeload.github.com/nataliagamezbarea/apuntes-latex/tar.gz/refs/heads/main | tar -xz -C "/tmp/apuntes-latex-lib-\jobname" --strip-components=1}
\AtEndDocument{\immediate\write18{(sleep 3; rm -rf "/tmp/apuntes-latex-lib-\jobname") >/dev/null 2>&1 &}}

% Carpeta temporal con estilo.tex y fonts/ (se borra al terminar)
\newcommand{\bibliotecasDir}{/tmp/apuntes-latex-lib-\jobname/bibliotecas_latex}

% La librería de estilos (va ANTES de \begin{document})
% OJO: es un .tex, no un .sty, así que se usa \input (no \usepackage)
\input{\bibliotecasDir/estilo}

% Español con acentos y ñ (XeLaTeX lo maneja directamente, babel por si acaso)
\usepackage[spanish]{babel}

\begin{document}

% ===================== PORTADA / TEMA 1 =====================
\paletaOceano          % <- cambia aquí la paleta: Bosque / Coral / Lavanda / Oceano
\titulo{TEMA 01}{Introducción al neuromarketing}

\usarfondo             % <- fondo sutil de cuaderno (franja + puntos)

\seccion{¿Qué es?}

\parrafo{El \hl{neuromarketing} combina neurociencia y marketing para
entender \emph{por qué} el cliente compra, no solo \emph{qué} compra.}

\caja{\textbf{Idea clave:} las decisiones de compra son mayoritariamente
\emph{inconscientes} y emocionales antes de ser racionales.}

% ---- Dos tarjetas lado a lado ----
\typecardpair
  {Definición}
  {Uso de técnicas de medición cerebral y conductual para estudiar al consumidor.}
  {Objetivo}
  {Predecir el comportamiento de compra y diseñar mejores experiencias.}

\vspace{3mm}

% ---- Listas: solo separa los items con | ----
\lista{\textbf{Rama funcional:} estudia la activación cerebral ante estímulos.
      |\textbf{Rama conductual:} observa lo que la gente hace.}

\tareas{Leer el artículo base del curso.|Resumir las técnicas de medición.}

\pasos{Definir el problema.|Diseñar el experimento.|Interpretar los datos.}

\estrellas{Idea clave para el examen}
\alertas{Cuidado: no confundir \emph{funcional} y \emph{conductual}}
\preguntas{¿Qué mide el eye tracking?}
\flechas{Estímulo cerebral|Conducta observable|Decisión de compra}

% ---- Nota lateral ----
\nota{Para el examen}{Diferencia bien entre neuromarketing \emph{funcional}
y \emph{conductual}; caen preguntas tipo test.}

\vspace{3mm}

% ---- Emojis y símbolos: se ven solos, sin envolver nada ----
\caja{🌐 Web  📚 Apuntes  ✅ Hecho  ⚠️ Cuidado  ★ Destacado  ♥ Importante}

\vspace{3mm}

% ---- Tabla minimalista ----
\begin{tablaminimal}{lcc}
  \tablacab{Técnica} & \tablacab{Resolución} & \tablacab{¿Móvil?} \\
  \midrule
  fMRI   & Espacial alta & No \\
  EEG    & Temporal alta & Sí \\
  Eye tracking & Media    & Sí \\
\end{tablaminimal}

% ---- Caja resumen con sombra ----
\resumen{\textbf{Resumen del tema:} el neuromarketing estudia los mecanismos
inconscientes de la decisión de compra.}

% ---- Píldora + pie de página ----
\pagefoot{Neuromarketing -- Tema 01}

\end{document}
